\section{Comparison between different fitting groups}
\label{app:group_comparison}
This appendix compares the fits produced by the MaCh3\cite{MaCh3:2018}, p-theta\cite{ptheta:2018} and VALOR fitting groups.

\subsection{Data fits}

\subsubsection{Results for \main}

The result for \main with reactor constraint and T2K only is shown in fig.~\ref{fig:dcp_vs_ssqth23_data_group_compare}.

\begin{figure*}[htb]
	\center
	\begin{subfigure}[b]{0.49\textwidth}
		\includegraphics[width=0.9\textwidth]{images/data/group_comparison/main_wrc_nh_data}
		\caption{\deltacp vs \ssqthtwothree NH with reactor constraint}
	\end{subfigure}
	\begin{subfigure}[b]{0.49\textwidth}
		\includegraphics[width=0.9\textwidth]{images/data/group_comparison/main_wrc_ih_data}
		\caption{\deltacp vs \ssqthtwothree IH with reactor constraint}
	\end{subfigure}
	\begin{subfigure}[b]{0.49\textwidth}
		\includegraphics[width=0.9\textwidth]{images/data/group_comparison/main_worc_nh_data}
		\caption{\deltacp vs \ssqthtwothree NH without reactor constraint}
	\end{subfigure}
	\begin{subfigure}[b]{0.49\textwidth}
		\includegraphics[width=0.9\textwidth]{images/data/group_comparison/main_worc_ih_data}
		\caption{\deltacp vs \ssqthtwothree IH without reactor constraint}
	\end{subfigure}
	\caption {
		Contours at 68\% and 90\% CL for \main with (top) and without (bottom) reactor constraint. Normal and inverted hierarchy contours are independent. All 5 samples were used to produce these fits.
	}
	\label{fig:dcp_vs_ssqth23_data_group_compare}
\end{figure*}

\clearpage

\subsubsection{Results for \appearance}
In \cref{fig:ssqth13_vs_dcp_data_group_compare} the result for \appearance with reactor constraint and T2K only is shown.

\begin{figure*}[htb]
	\center
	\begin{subfigure}[b]{0.49\textwidth}
		\includegraphics[width=0.9\textwidth]{images/data/group_comparison/appearance_wrc_nh_data}
		\caption{\appearance NH with reactor constraint}
	\end{subfigure}
	\begin{subfigure}[b]{0.49\textwidth}
		\includegraphics[width=0.9\textwidth]{images/data/group_comparison/appearance_wrc_ih_data}
		\caption{\appearance IH with reactor constraint}
	\end{subfigure}
	\begin{subfigure}[b]{0.49\textwidth}
		\includegraphics[width=0.9\textwidth]{images/data/group_comparison/appearance_worc_nh_data}
		\caption{\appearance NH without reactor constraint}
	\end{subfigure}
	\begin{subfigure}[b]{0.49\textwidth}
		\includegraphics[width=0.9\textwidth]{images/data/group_comparison/appearance_worc_ih_data}
		\caption{\appearance IH without reactor constraint}
	\end{subfigure}
	\caption {
		Contours at 68\% and 90\% CL for \appearance with and without reactor constraint. Normal and inverted hierarchy contours are independent. All 5 samples were used to produce these fits.
	}
	\label{fig:ssqth13_vs_dcp_data_group_compare}
\end{figure*}
\clearpage

\subsubsection{Results for \disappearance}

The result for \disappearance with reactor constraint and T2K only is shown in fig.~\ref{fig:ssqth23_vs_dmsq_data_group_compare}. The difference between p-theta and the other fitting groups has been seen previously and was driven by the difference in variables and binning between the fitters.

\begin{figure*}[htb]
	\center
	\begin{subfigure}[b]{0.49\textwidth}
		\includegraphics[width=0.9\textwidth]{images/data/group_comparison/disappearance_wrc_nh_data}
		\caption{\disappearance NH with reactor constraint}
	\end{subfigure}
	\begin{subfigure}[b]{0.49\textwidth}
		\includegraphics[width=0.9\textwidth]{images/data/group_comparison/disappearance_wrc_ih_data}
		\caption{\disappearance IH with reactor constraint}
	\end{subfigure}
	\begin{subfigure}[b]{0.49\textwidth}
		\includegraphics[width=0.9\textwidth]{images/data/group_comparison/disappearance_worc_nh_data}
		\caption{\disappearance NH without reactor constraint}
	\end{subfigure}
	\begin{subfigure}[b]{0.49\textwidth}
		\includegraphics[width=0.9\textwidth]{images/data/group_comparison/disappearance_worc_ih_data}
		\caption{\disappearance IH without reactor constraint}
	\end{subfigure}
	\begin{subfigure}[b]{0.49\textwidth}
%		\includegraphics[width=0.9\textwidth]{images/data/group_comparison/disappearance_wrc_nh_ptheta_erec}
		\caption{\disappearance NH with reactor constraint}
	\end{subfigure}
	\begin{subfigure}[b]{0.49\textwidth}
%		\includegraphics[width=0.9\textwidth]{images/data/group_comparison/disappearance_wrc_ih_ptheta_erec}
		\caption{\disappearance IH with reactor constraint}
	\end{subfigure}
	\caption {
		Contours at 68\% and 90\% CL for \disappearance with (top) and without (middle) reactor constraint and comparison to p-theta reactor fit with $E_r-\theta$ templates (bottom). Normal and inverted hierarchy contours are independent. All 5 samples were used to produce these fits.
	}
	\label{fig:ssqth23_vs_dmsq_data_group_compare}
\end{figure*}

\clearpage

\subsubsection{Results for \deltacp}

The result for \deltacp with reactor constraint and T2K only is shown in fig.~\ref{fig:dcp_data_group_compare}.

\begin{figure*}[htb]
	\center
	\begin{subfigure}[b]{0.49\textwidth}
		\includegraphics[width=0.9\textwidth]{images/data/group_comparison/deltacp_wrc_nh_data}
		\caption{\deltacp NH with reactor constraint}
	\end{subfigure}
	\begin{subfigure}[b]{0.49\textwidth}
		\includegraphics[width=0.9\textwidth]{images/data/group_comparison/deltacp_wrc_ih_data}
		\caption{\deltacp IH with reactor constraint}
	\end{subfigure}
	\begin{subfigure}[b]{0.49\textwidth}
		\includegraphics[width=0.9\textwidth]{images/data/group_comparison/deltacp_worc_nh_data}
		\caption{\deltacp NH without reactor constraint}
	\end{subfigure}
	\begin{subfigure}[b]{0.49\textwidth}
		\includegraphics[width=0.9\textwidth]{images/data/group_comparison/deltacp_worc_ih_data}
		\caption{\deltacp IH without reactor constraint}
	\end{subfigure}
	\caption {
		The expected $\Delta\chi^2$ distribution as a function of \deltacp with and without reactor constraint. Normal and inverted hierarchy contours are independent. All 5 samples were used to produce these fits.
	}
	\label{fig:dcp_data_group_compare}
\end{figure*}
\clearpage

\subsubsection{Results for \ssqthonethree}

The result for \ssqthonethree (without reactor constraint) is shown in fig.~\ref{fig:th13_data_group_compare}.

\begin{figure*}[htb]
	\center
	\begin{subfigure}[b]{0.49\textwidth}
		\includegraphics[width=0.9\textwidth]{images/data/group_comparison/theta13_worc_nh_data}
		\caption{\ssqthonethree NH without reactor constraint}
	\end{subfigure}
	\begin{subfigure}[b]{0.49\textwidth}
		\includegraphics[width=0.9\textwidth]{images/data/group_comparison/theta13_worc_ih_data}
		\caption{\ssqthonethree IH without reactor constraint}
	\end{subfigure}
	\caption {
		The expected $\Delta\chi^2$ distribution as a function of \ssqthonethree without reactor constraint. Normal and inverted hierarchy contours are independent. All 5 samples were used to produce these fits.
	}
	\label{fig:th13_data_group_compare}
\end{figure*}

\clearpage

\subsubsection{Results for \ssqthtwothree}

The result for \ssqthtwothree with reactor constraint and T2K only is shown in fig.~\ref{fig:th23_data_group_compare}.

\begin{figure*}[htb]
	\center
	\begin{subfigure}[b]{0.49\textwidth}
		\includegraphics[width=0.9\textwidth]{images/data/group_comparison/theta23_wrc_nh_data}
		\caption{\ssqthtwothree NH with reactor constraint}
	\end{subfigure}
	\begin{subfigure}[b]{0.49\textwidth}
		\includegraphics[width=0.9\textwidth]{images/data/group_comparison/theta23_wrc_ih_data}
		\caption{\ssqthtwothree IH with reactor constraint}
	\end{subfigure}
	\begin{subfigure}[b]{0.49\textwidth}
		\includegraphics[width=0.9\textwidth]{images/data/group_comparison/theta23_worc_nh_data}
		\caption{\ssqthtwothree NH without reactor constraint}
	\end{subfigure}
	\begin{subfigure}[b]{0.49\textwidth}
		\includegraphics[width=0.9\textwidth]{images/data/group_comparison/theta23_worc_ih_data}
		\caption{\ssqthtwothree IH without reactor constraint}
	\end{subfigure}
	\caption {
		The expected $\Delta\chi^2$ distribution as a function of \ssqthtwothree with and without reactor constraint. Normal and inverted hierarchy contours are independent. All 5 samples were used to produce these fits.
	}
	\label{fig:th23_data_group_compare}
\end{figure*}
\clearpage

\subsubsection{Results for \dmsq}

The result for \dmsq with reactor constraint and T2K only is shown in fig.~\ref{fig:dmsq_data_group_compare}.

\begin{figure*}[htb]
	\center
	\begin{subfigure}[b]{0.49\textwidth}
		\includegraphics[width=0.9\textwidth]{images/data/group_comparison/dmsq_wrc_nh_data}
		\caption{\dmsq NH with reactor constraint}
	\end{subfigure}
	\begin{subfigure}[b]{0.49\textwidth}
		\includegraphics[width=0.9\textwidth]{images/data/group_comparison/dmsq_wrc_ih_data}
		\caption{\dmsq IH with reactor constraint}
	\end{subfigure}
	\begin{subfigure}[b]{0.49\textwidth}
		\includegraphics[width=0.9\textwidth]{images/data/group_comparison/dmsq_worc_nh_data}
		\caption{\dmsq NH without reactor constraint}
	\end{subfigure}
	\begin{subfigure}[b]{0.49\textwidth}
		\includegraphics[width=0.9\textwidth]{images/data/group_comparison/dmsq_worc_ih_data}
		\caption{\dmsq IH without reactor constraint}
	\end{subfigure}
	\caption {
		The expected $\Delta\chi^2$ distribution as a function of \dmsqtwothree with and without reactor constraint. Normal and inverted hierarchy contours are independent. All 5 samples were used to produce these fits.
	}
	\label{fig:dmsq_data_group_compare}
\end{figure*}
\FloatBarrier
\clearpage

\subsection{Asimov A fits}

\subsubsection{Results for \appearance}
In \cref{fig:ssqth13_vs_dcp_asimova_group_compare} the result for \appearance with reactor constraint and T2K only is shown.

\begin{figure*}[htb]
	\center
	\begin{subfigure}[b]{0.49\textwidth}
		\includegraphics[width=0.9\textwidth]{images/sensitivity/group_comparison/appearance_wrc_nh_asimova}
		\caption{\appearance NH with reactor constraint}
	\end{subfigure}
	\begin{subfigure}[b]{0.49\textwidth}
		\includegraphics[width=0.9\textwidth]{images/sensitivity/group_comparison/appearance_wrc_ih_asimova}
		\caption{\appearance IH with reactor constraint}
	\end{subfigure}
	\begin{subfigure}[b]{0.49\textwidth}
		\includegraphics[width=0.9\textwidth]{images/sensitivity/group_comparison/appearance_worc_nh_asimova}
		\caption{\appearance NH without reactor constraint}
	\end{subfigure}
	\begin{subfigure}[b]{0.49\textwidth}
		\includegraphics[width=0.9\textwidth]{images/sensitivity/group_comparison/appearance_worc_ih_asimova}
		\caption{\appearance IH without reactor constraint}
	\end{subfigure}
	\caption {
		Contours at 68\% and 90\% CL for \appearance with and without reactor constraint. Normal and inverted hierarchy contours are independent. All 5 samples were used to produce these fits.
	}
	\label{fig:ssqth13_vs_dcp_asimova_group_compare}
\end{figure*}
\clearpage

\subsubsection{Results for \disappearance}

The result for \disappearance with reactor constraint and T2K only is shown in fig.~\ref{fig:ssqth23_vs_dmsq_asimova_group_compare}.

\begin{figure*}[htb]
	\center
	\begin{subfigure}[b]{0.49\textwidth}
		\includegraphics[width=0.9\textwidth]{images/sensitivity/group_comparison/disappearance_wrc_nh_asimova}
		\caption{\disappearance NH with reactor constraint}
	\end{subfigure}
	\begin{subfigure}[b]{0.49\textwidth}
		\includegraphics[width=0.9\textwidth]{images/sensitivity/group_comparison/disappearance_wrc_ih_asimova}
		\caption{\disappearance IH with reactor constraint}
	\end{subfigure}
	\begin{subfigure}[b]{0.49\textwidth}
		\includegraphics[width=0.9\textwidth]{images/sensitivity/group_comparison/disappearance_worc_nh_asimova}
		\caption{\disappearance NH without reactor constraint}
	\end{subfigure}
	\begin{subfigure}[b]{0.49\textwidth}
		\includegraphics[width=0.9\textwidth]{images/sensitivity/group_comparison/disappearance_worc_ih_asimova}
		\caption{\disappearance IH without reactor constraint}
	\end{subfigure}
	\caption {
		Contours at 68\% and 90\% CL for \disappearance with and without reactor constraint. Normal and inverted hierarchy contours are independent. All 5 samples were used to produce these fits.
	}
	\label{fig:ssqth23_vs_dmsq_asimova_group_compare}
\end{figure*}
\clearpage

\subsubsection{Results for \deltacp}

The result for \deltacp with reactor constraint and T2K only is shown in fig.~\ref{fig:dcp_asimova_group_compare}.

\begin{figure*}[htb]
	\center
	\begin{subfigure}[b]{0.49\textwidth}
		\includegraphics[width=0.9\textwidth]{images/sensitivity/group_comparison/deltacp_wrc_nh_asimova}
		\caption{\deltacp NH with reactor constraint}
	\end{subfigure}
	\begin{subfigure}[b]{0.49\textwidth}
		\includegraphics[width=0.9\textwidth]{images/sensitivity/group_comparison/deltacp_wrc_ih_asimova}
		\caption{\deltacp IH with reactor constraint}
	\end{subfigure}
	\begin{subfigure}[b]{0.49\textwidth}
		\includegraphics[width=0.9\textwidth]{images/sensitivity/group_comparison/deltacp_worc_nh_asimova}
		\caption{\deltacp NH without reactor constraint}
	\end{subfigure}
	\begin{subfigure}[b]{0.49\textwidth}
		\includegraphics[width=0.9\textwidth]{images/sensitivity/group_comparison/deltacp_worc_ih_asimova}
		\caption{\deltacp IH without reactor constraint}
	\end{subfigure}
	\caption {
		The expected $\Delta\chi^2$ distribution as a function of \deltacp with and without reactor constraint. Normal and inverted hierarchy contours are independent. All 5 samples were used to produce these fits.
	}
	\label{fig:dcp_asimova_group_compare}
\end{figure*}
\clearpage

\subsubsection{Results for \ssqthonethree}

The result for \ssqthonethree (without reactor constraint) is shown in fig.~\ref{fig:th13_asimova_group_compare}.

\begin{figure*}[htb]
	\center
	\begin{subfigure}[b]{0.49\textwidth}
		\includegraphics[width=0.9\textwidth]{images/sensitivity/group_comparison/theta13_worc_nh_asimova}
		\caption{\ssqthonethree NH without reactor constraint}
	\end{subfigure}
	\begin{subfigure}[b]{0.49\textwidth}
		\includegraphics[width=0.9\textwidth]{images/sensitivity/group_comparison/theta13_worc_ih_asimova}
		\caption{\ssqthonethree IH without reactor constraint}
	\end{subfigure}
	\caption {
		The expected $\Delta\chi^2$ distribution as a function of \ssqthonethree without reactor constraint. Normal and inverted hierarchy contours are independent. All 5 samples were used to produce these fits.
	}
	\label{fig:th13_asimova_group_compare}
\end{figure*}

\clearpage

\subsubsection{Results for \ssqthtwothree}

The result for \ssqthtwothree with reactor constraint and T2K only is shown in fig.~\ref{fig:th23_asimova_group_compare}.

\begin{figure*}[htb]
	\center
	\begin{subfigure}[b]{0.49\textwidth}
		\includegraphics[width=0.9\textwidth]{images/sensitivity/group_comparison/theta23_wrc_nh_asimova}
		\caption{\ssqthtwothree NH with reactor constraint}
	\end{subfigure}
	\begin{subfigure}[b]{0.49\textwidth}
		\includegraphics[width=0.9\textwidth]{images/sensitivity/group_comparison/theta23_wrc_ih_asimova}
		\caption{\ssqthtwothree IH with reactor constraint}
	\end{subfigure}
	\begin{subfigure}[b]{0.49\textwidth}
		\includegraphics[width=0.9\textwidth]{images/sensitivity/group_comparison/theta23_worc_nh_asimova}
		\caption{\ssqthtwothree NH without reactor constraint}
	\end{subfigure}
	\begin{subfigure}[b]{0.49\textwidth}
		\includegraphics[width=0.9\textwidth]{images/sensitivity/group_comparison/theta23_worc_ih_asimova}
		\caption{\ssqthtwothree IH without reactor constraint}
	\end{subfigure}
	\caption {
		The expected $\Delta\chi^2$ distribution as a function of \ssqthtwothree with and without reactor constraint. Normal and inverted hierarchy contours are independent. All 5 samples were used to produce these fits.
	}
	\label{fig:th23_asimova_group_compare}
\end{figure*}
\clearpage

\subsubsection{Results for \dmsq}

The result for \dmsq with reactor constraint and T2K only is shown in fig.~\ref{fig:dmsq_asimova_group_compare}.

\begin{figure*}[htb]
	\center
	\begin{subfigure}[b]{0.49\textwidth}
		\includegraphics[width=0.9\textwidth]{images/sensitivity/group_comparison/dmsq_wrc_nh_asimova}
		\caption{\dmsq NH with reactor constraint}
	\end{subfigure}
	\begin{subfigure}[b]{0.49\textwidth}
		\includegraphics[width=0.9\textwidth]{images/sensitivity/group_comparison/dmsq_wrc_ih_asimova}
		\caption{\dmsq IH with reactor constraint}
	\end{subfigure}
	\begin{subfigure}[b]{0.49\textwidth}
		\includegraphics[width=0.9\textwidth]{images/sensitivity/group_comparison/dmsq_worc_nh_asimova}
		\caption{\dmsq NH without reactor constraint}
	\end{subfigure}
	\begin{subfigure}[b]{0.49\textwidth}
		\includegraphics[width=0.9\textwidth]{images/sensitivity/group_comparison/dmsq_worc_ih_asimova}
		\caption{\dmsq IH without reactor constraint}
	\end{subfigure}
	\caption {
		The expected $\Delta\chi^2$ distribution as a function of \dmsqtwothree with and without reactor constraint. Normal and inverted hierarchy contours are independent. All 5 samples were used to produce these fits.
	}
	\label{fig:dmsq_asimova_group_compare}
\end{figure*}
\FloatBarrier
\clearpage

\subsection{Asimov B fits}

\subsubsection{Results for \appearance}
In \cref{fig:ssqth13_vs_dcp_asimovb_group_compare} the result for \appearance with reactor constraint and T2K only is shown.

\begin{figure*}[htb]
	\center
	\begin{subfigure}[b]{0.49\textwidth}
		\includegraphics[width=0.9\textwidth]{images/sensitivity/group_comparison/appearance_wrc_nh_asimovb}
		\caption{\appearance NH with reactor constraint}
	\end{subfigure}
	\begin{subfigure}[b]{0.49\textwidth}
		\includegraphics[width=0.9\textwidth]{images/sensitivity/group_comparison/appearance_wrc_ih_asimovb}
		\caption{\appearance IH with reactor constraint}
	\end{subfigure}
	\begin{subfigure}[b]{0.49\textwidth}
		\includegraphics[width=0.9\textwidth]{images/sensitivity/group_comparison/appearance_worc_nh_asimovb}
		\caption{\appearance NH without reactor constraint}
	\end{subfigure}
	\begin{subfigure}[b]{0.49\textwidth}
		\includegraphics[width=0.9\textwidth]{images/sensitivity/group_comparison/appearance_worc_ih_asimovb}
		\caption{\appearance IH without reactor constraint}
	\end{subfigure}
	\caption {
		Contours at 68\% and 90\% CL for \appearance with and without reactor constraint. Normal and inverted hierarchy contours are independent. All 5 samples were used to produce these fits.
	}
	\label{fig:ssqth13_vs_dcp_asimovb_group_compare}
\end{figure*}
\clearpage

\subsubsection{Results for \disappearance}

The result for \disappearance with reactor constraint and T2K only is shown in fig.~\ref{fig:ssqth23_vs_dmsq_asimovb_group_compare}.

\begin{figure*}[htb]
	\center
	\begin{subfigure}[b]{0.49\textwidth}
		\includegraphics[width=0.9\textwidth]{images/sensitivity/group_comparison/disappearance_wrc_nh_asimovb}
		\caption{\disappearance NH with reactor constraint}
	\end{subfigure}
	\begin{subfigure}[b]{0.49\textwidth}
		\includegraphics[width=0.9\textwidth]{images/sensitivity/group_comparison/disappearance_wrc_ih_asimovb}
		\caption{\disappearance IH with reactor constraint}
	\end{subfigure}
	\begin{subfigure}[b]{0.49\textwidth}
		\includegraphics[width=0.9\textwidth]{images/sensitivity/group_comparison/disappearance_worc_nh_asimovb}
		\caption{\disappearance NH without reactor constraint}
	\end{subfigure}
	\begin{subfigure}[b]{0.49\textwidth}
		\includegraphics[width=0.9\textwidth]{images/sensitivity/group_comparison/disappearance_worc_ih_asimovb}
		\caption{\disappearance IH without reactor constraint}
	\end{subfigure}
	\caption {
		Contours at 68\% and 90\% CL for \disappearance with and without reactor constraint. Normal and inverted hierarchy contours are independent. All 5 samples were used to produce these fits.
	}
	\label{fig:ssqth23_vs_dmsq_asimovb_group_compare}
\end{figure*}
\clearpage

\subsubsection{Results for \deltacp}

The result for \deltacp with reactor constraint and T2K only is shown in fig.~\ref{fig:dcp_asimovb_group_compare}.

\begin{figure*}[htb]
	\center
	\begin{subfigure}[b]{0.49\textwidth}
		\includegraphics[width=0.9\textwidth]{images/sensitivity/group_comparison/deltacp_wrc_nh_asimovb}
		\caption{\deltacp NH with reactor constraint}
	\end{subfigure}
	\begin{subfigure}[b]{0.49\textwidth}
		\includegraphics[width=0.9\textwidth]{images/sensitivity/group_comparison/deltacp_wrc_ih_asimovb}
		\caption{\deltacp IH with reactor constraint}
	\end{subfigure}
	\begin{subfigure}[b]{0.49\textwidth}
		\includegraphics[width=0.9\textwidth]{images/sensitivity/group_comparison/deltacp_worc_nh_asimovb}
		\caption{\deltacp NH without reactor constraint}
	\end{subfigure}
	\begin{subfigure}[b]{0.49\textwidth}
		\includegraphics[width=0.9\textwidth]{images/sensitivity/group_comparison/deltacp_worc_ih_asimovb}
		\caption{\deltacp IH without reactor constraint}
	\end{subfigure}
	\caption {
		The expected $\Delta\chi^2$ distribution as a function of \deltacp with and without reactor constraint. Normal and inverted hierarchy contours are independent. All 5 samples were used to produce these fits.
	}
	\label{fig:dcp_asimovb_group_compare}
\end{figure*}
\clearpage

\subsubsection{Results for \ssqthonethree}

The result for \ssqthonethree (without reactor constraint) is shown in fig.~\ref{fig:th13_asimovb_group_compare}.

\begin{figure*}[htb]
	\center
	\begin{subfigure}[b]{0.49\textwidth}
		\includegraphics[width=0.9\textwidth]{images/sensitivity/group_comparison/theta13_worc_nh_asimovb}
		\caption{\ssqthonethree NH without reactor constraint}
	\end{subfigure}
	\begin{subfigure}[b]{0.49\textwidth}
		\includegraphics[width=0.9\textwidth]{images/sensitivity/group_comparison/theta13_worc_ih_asimovb}
		\caption{\ssqthonethree IH without reactor constraint}
	\end{subfigure}
	\caption {
		The expected $\Delta\chi^2$ distribution as a function of \ssqthonethree without reactor constraint. Normal and inverted hierarchy contours are independent. All 5 samples were used to produce these fits.
	}
	\label{fig:th13_asimovb_group_compare}
\end{figure*}

\clearpage

\subsubsection{Results for \ssqthtwothree}

The result for \ssqthtwothree with reactor constraint and T2K only is shown in fig.~\ref{fig:th23_asimovb_group_compare}.

\begin{figure*}[htb]
	\center
	\begin{subfigure}[b]{0.49\textwidth}
		\includegraphics[width=0.9\textwidth]{images/sensitivity/group_comparison/theta23_wrc_nh_asimovb}
		\caption{\ssqthtwothree NH with reactor constraint}
	\end{subfigure}
	\begin{subfigure}[b]{0.49\textwidth}
		\includegraphics[width=0.9\textwidth]{images/sensitivity/group_comparison/theta23_wrc_ih_asimovb}
		\caption{\ssqthtwothree IH with reactor constraint}
	\end{subfigure}
	\begin{subfigure}[b]{0.49\textwidth}
		\includegraphics[width=0.9\textwidth]{images/sensitivity/group_comparison/theta23_worc_nh_asimovb}
		\caption{\ssqthtwothree NH without reactor constraint}
	\end{subfigure}
	\begin{subfigure}[b]{0.49\textwidth}
		\includegraphics[width=0.9\textwidth]{images/sensitivity/group_comparison/theta23_worc_ih_asimovb}
		\caption{\ssqthtwothree IH without reactor constraint}
	\end{subfigure}
	\caption {
		The expected $\Delta\chi^2$ distribution as a function of \ssqthtwothree with and without reactor constraint. Normal and inverted hierarchy contours are independent. All 5 samples were used to produce these fits.
	}
	\label{fig:th23_asimovb_group_compare}
\end{figure*}
\clearpage

\subsubsection{Results for \dmsq}

The result for \dmsq with reactor constraint and T2K only is shown in fig.~\ref{fig:dmsq_asimovb_group_compare}.

\begin{figure*}[htb]
	\center
	\begin{subfigure}[b]{0.49\textwidth}
		\includegraphics[width=0.9\textwidth]{images/sensitivity/group_comparison/dmsq_wrc_nh_asimovb}
		\caption{\dmsq NH with reactor constraint}
	\end{subfigure}
	\begin{subfigure}[b]{0.49\textwidth}
		\includegraphics[width=0.9\textwidth]{images/sensitivity/group_comparison/dmsq_wrc_ih_asimovb}
		\caption{\dmsq IH with reactor constraint}
	\end{subfigure}
	\begin{subfigure}[b]{0.49\textwidth}
		\includegraphics[width=0.9\textwidth]{images/sensitivity/group_comparison/dmsq_worc_nh_asimovb}
		\caption{\dmsq NH without reactor constraint}
	\end{subfigure}
	\begin{subfigure}[b]{0.49\textwidth}
		\includegraphics[width=0.9\textwidth]{images/sensitivity/group_comparison/dmsq_worc_ih_asimovb}
		\caption{\dmsq IH without reactor constraint}
	\end{subfigure}
	\caption {
		The expected $\Delta\chi^2$ distribution as a function of \dmsqtwothree with and without reactor constraint. Normal and inverted hierarchy contours are independent. All 5 samples were used to produce these fits.
	}
	\label{fig:dmsq_asimovb_group_compare}
\end{figure*}
\FloatBarrier
\clearpage